% !TEX root = ../main.tex
\chapter{Probability Distributions}
\label{ch:distributions}

Probability distributions are the mathematical objects of
Bayesian inference, generative modeling, variational autoencoders,
and reinforcement learning. \Lucid\ exposes a comprehensive
distributions subpackage at \code{lucid.distributions} that
provides parameterised distribution objects with sample, density,
log-density, KL divergence, and entropy methods.

This chapter describes the subpackage's design, the supported
distributions, the reparameterisation trick for differentiable
sampling, and the KL-divergence registry.

\section{Overview}
\label{sec:distributions:overview}

The subpackage follows the design of \refframework's
\code{torch.distributions} and TensorFlow Probability: a
\code{Distribution} abstract base, parameterised concrete
subclasses, and a fluent interface for sample / log\_prob / mean
/ variance / entropy / KL.

\subsection{The 33 distributions}
\label{ssec:distributions:catalogue}

\Lucid\ supports 33 distributions across four categories:

\begin{itemize}
  \item \textbf{Continuous}: \code{Normal}, \code{Laplace},
    \code{Cauchy}, \code{StudentT}, \code{Gumbel},
    \code{Exponential}, \code{Gamma}, \code{Beta},
    \code{Dirichlet}, \code{LogNormal}, \code{Weibull},
    \code{Pareto}, \code{HalfNormal}, \code{HalfCauchy},
    \code{Uniform}, \code{VonMises}.
  \item \textbf{Discrete}: \code{Bernoulli}, \code{Binomial},
    \code{Categorical}, \code{OneHotCategorical}, \code{Poisson},
    \code{Geometric}, \code{NegativeBinomial},
    \code{Multinomial}.
  \item \textbf{Multivariate}: \code{MultivariateNormal},
    \code{Wishart}, \code{LKJCholesky}, \code{Dirichlet}.
  \item \textbf{Relaxed / continuous discrete}:
    \code{RelaxedBernoulli}, \code{RelaxedOneHotCategorical}.
  \item \textbf{Transformed / mixture}: \code{TransformedDistribution},
    \code{MixtureSameFamily}, \code{Independent}.
\end{itemize}

\section{The Distribution Interface}
\label{sec:distributions:interface}

Every concrete distribution implements:

\begin{itemize}
  \item \code{sample(shape=())}: draw a sample of the requested
    batch shape.
  \item \code{rsample(shape=())}: differentiable sample (for
    reparameterisable distributions).
  \item \code{log\_prob(value)}: log probability density / mass at
    \code{value}.
  \item \code{cdf(value)}: cumulative distribution function (for
    those that have closed-form CDF).
  \item \code{icdf(value)}: inverse CDF.
  \item \code{mean, variance, stddev, entropy}: properties /
    methods returning the standard statistics.
\end{itemize}

\subsection{Shape conventions}
\label{ssec:distributions:shapes}

A distribution has three shape conceptions:

\begin{itemize}
  \item \textbf{batch\_shape}: independent distributions arranged
    in a batch. For a Normal with mean shape $(B,)$ and stddev
    shape $(B,)$, the batch shape is $(B,)$.
  \item \textbf{event\_shape}: the shape of a single draw. For a
    scalar Normal, event\_shape is $()$. For
    \code{MultivariateNormal} of dimension $D$, event\_shape is
    $(D,)$.
  \item \textbf{sample\_shape}: requested shape of the sample;
    prepended to batch\_shape and event\_shape.
\end{itemize}

\code{sample(shape=(K,))} of a \code{Normal} with batch shape
$(B,)$ returns shape $(K, B)$. \code{sample(shape=(K, M))} of
\code{MultivariateNormal} with batch shape $()$ and event shape
$(D,)$ returns shape $(K, M, D)$.

\subsection{The Independent wrapper}
\label{ssec:distributions:independent}

\code{Independent(base, n)} reinterprets the last $n$ dimensions
of \code{base}'s batch\_shape as event\_shape. Useful when the
user has a vector of independent scalar Normals and wants
\code{log\_prob} to sum across them rather than return a
per-element tensor:

\begin{lstlisting}[style=lucid,language=LucidPython]
import lucid
from lucid.distributions import Normal, Independent

# Batch of D-dim diagonal Normals
mu = lucid.zeros(B, D)
sigma = lucid.ones(B, D)
base = Normal(mu, sigma)             # batch (B, D), event ()
dist = Independent(base, 1)          # batch (B,), event (D,)

x = dist.sample()                    # shape (B, D)
lp = dist.log_prob(x)                # shape (B,), summed over D
\end{lstlisting}

\section{The Reparameterisation Trick}
\label{sec:distributions:reparam}

For training a model whose loss includes an expectation over a
distribution (e.g., VAE evidence lower bound), the gradient of
the expectation with respect to the distribution's parameters
needs to be computable. The reparameterisation trick is the
standard mechanism.

\subsection{The trick}
\label{ssec:distributions:reparam_trick}

Suppose a sample $x \sim p_\theta$ can be written as $x =
g_\theta(\epsilon)$ where $\epsilon$ comes from a fixed,
parameter-free distribution $q$. Then
\[
\nabla_\theta \,\mathbb{E}_{x \sim p_\theta}[f(x)]
\;=\; \nabla_\theta \,\mathbb{E}_{\epsilon \sim q}[f(g_\theta(\epsilon))]
\;=\; \mathbb{E}_{\epsilon \sim q}\!\left[\nabla_\theta f(g_\theta(\epsilon))\right].
\]
The expectation and gradient commute because $q$ does not depend
on $\theta$. The gradient passes through the deterministic
$g_\theta$; the stochastic $\epsilon$ is a constant from autograd's
view.

For Gaussian:
\[
x \sim \mathcal{N}(\mu, \sigma^2)
\;\Longleftrightarrow\;
x = \mu + \sigma \cdot \epsilon,
\quad \epsilon \sim \mathcal{N}(0, 1).
\]
The gradients are then immediate:
$\partial x / \partial \mu = 1$ and $\partial x / \partial \sigma
= \epsilon$. Autograd handles both automatically.

For other reparameterisable distributions the transform is
analogous: $\mathrm{Exponential}(\lambda)$ as $x = -\log(u)/\lambda$
with $u \sim \mathcal{U}(0, 1)$; $\mathrm{Uniform}(a, b)$ as $x =
a + (b - a) u$; etc.

\subsection{rsample}
\label{ssec:distributions:rsample}

\code{rsample()} returns a reparameterised sample for
distributions that admit the trick. The implementation calls
\code{torch.randn}-style noise and applies the differentiable
transformation.

Distributions that support \code{rsample}: \code{Normal},
\code{Uniform}, \code{Laplace}, \code{Gumbel}, \code{Exponential},
\code{Pareto}, \code{LogNormal}, \code{StudentT},
\code{HalfNormal}, \code{HalfCauchy}, \code{Cauchy},
\code{Beta}, \code{Dirichlet}, \code{Gamma},
\code{MultivariateNormal}.

\code{sample} (non-r) is the version that produces a sample
without the reparameterisation graph; useful when gradients are
not needed.

\subsection{Score-function estimator}
\label{ssec:distributions:score_function}

For non-reparameterisable distributions (discrete distributions,
some heavy-tailed continuous), the score-function estimator
(REINFORCE) is the alternative:
\[
\nabla_\theta \mathbb{E}_{x \sim p_\theta}[f(x)] =
\mathbb{E}_{x \sim p_\theta}[f(x) \nabla_\theta \log p_\theta(x)].
\]

\Lucid\ does not provide a built-in score-function gradient
helper, but the user can implement it explicitly:

\begin{lstlisting}[style=lucid,language=LucidPython]
x = dist.sample()                # non-differentiable sample
log_p = dist.log_prob(x)         # differentiable in dist's params
loss = -(reward * log_p).mean()  # score-function objective
\end{lstlisting}

The pattern is standard for policy gradient in RL.

\section{KL Divergence}
\label{sec:distributions:kl}

The Kullback--Leibler divergence
\[
D_{\mathrm{KL}}(p \,\|\, q)
\;=\; \mathbb{E}_{x \sim p}\!\left[\log \frac{p(x)}{q(x)}\right]
\;=\; \int p(x) \log \frac{p(x)}{q(x)} \, dx
\]
measures the directed distance from $p$ to $q$. It is
non-negative, zero iff $p = q$ almost everywhere, asymmetric in
its arguments, and unbounded above. It is the standard objective
in variational inference, the regulariser in $\beta$-VAE, and the
core quantity in information theory's notion of inefficient
coding.

For two univariate Gaussians the closed form is
\[
D_{\mathrm{KL}}\!\bigl(\mathcal{N}(\mu_1, \sigma_1^2) \,\|\, \mathcal{N}(\mu_2, \sigma_2^2)\bigr)
\;=\; \log \frac{\sigma_2}{\sigma_1}
   + \frac{\sigma_1^2 + (\mu_1 - \mu_2)^2}{2\sigma_2^2}
   - \frac{1}{2}.
\]
For two categorical distributions over the same support,
\[
D_{\mathrm{KL}}(p \,\|\, q) \;=\; \sum_i p_i \log \frac{p_i}{q_i}.
\]
The first form appears in every VAE loss; the second in
distillation, label-smoothing, and policy-gradient
regularisation.

\subsection{Closed-form KLs}
\label{ssec:distributions:closed_form_kl}

For many pairs $(p, q)$ of common distributions, the KL has a
closed-form expression that avoids the expensive Monte-Carlo
estimate. \Lucid\ maintains a registry of these in
\code{lucid/distributions/kl.py}:

\begin{itemize}
  \item Normal -- Normal: closed form involving mean diff and
    log-stddev ratio.
  \item Categorical -- Categorical: $\sum_i p_i \log(p_i / q_i)$.
  \item Bernoulli -- Bernoulli: closed form.
  \item Gamma -- Gamma: closed form involving lgamma and digamma.
  \item Beta -- Beta: closed form.
  \item Dirichlet -- Dirichlet: closed form.
  \item MultivariateNormal -- MultivariateNormal: closed form
    involving log-det and trace.
\end{itemize}

The registry has 20 pairs as of \Lucid\ 3.5.

\subsection{Monte-Carlo fallback}
\label{ssec:distributions:mc_kl}

For pairs not in the registry, \code{kl\_divergence(p, q)} falls
back to an MC estimate:
\[
\widehat{D_{\text{KL}}}(p \| q) = \frac{1}{N} \sum_{i=1}^N \log
\frac{p(x_i)}{q(x_i)}, \quad x_i \sim p.
\]

The MC estimate is noisy; the user should average over many
samples for a meaningful loss.

\subsection{Registering a custom KL}
\label{ssec:distributions:custom_kl}

The user can register a custom closed-form KL via a decorator:

\begin{lstlisting}[style=lucid,language=LucidPython]
from lucid.distributions import register_kl, MyDist

@register_kl(MyDist, MyDist)
def _kl_mydist_mydist(p, q):
    return ...  # closed-form expression
\end{lstlisting}

\code{kl\_divergence(p, q)} will then pick the registered
implementation.

\section{Worked Example: VAE Loss}
\label{sec:distributions:vae_example}

The Variational Autoencoder's loss is the evidence lower bound
(ELBO):
\[
\mathcal{L} = \mathbb{E}_{q(z|x)}[\log p(x|z)] -
D_{\text{KL}}(q(z|x) \| p(z)).
\]

In \Lucid:

\begin{lstlisting}[style=lucid,language=LucidPython]
import lucid
from lucid.distributions import Normal, kl_divergence

# Encoder produces mean and log-stddev of q(z|x)
mu, log_sigma = encoder(x)
sigma = log_sigma.exp()
q_z = Normal(mu, sigma)

# Sample z via reparameterisation
z = q_z.rsample()

# Decoder produces parameters of p(x|z)
x_recon = decoder(z)
p_x_given_z = Normal(x_recon, lucid.tensor(0.1))
recon_loss = -p_x_given_z.log_prob(x).sum(dim=-1).mean()

# KL divergence to the prior p(z) = N(0, I)
prior = Normal(lucid.zeros_like(mu), lucid.ones_like(sigma))
kl_loss = kl_divergence(q_z, prior).sum(dim=-1).mean()

loss = recon_loss + kl_loss
loss.backward()
\end{lstlisting}

The \code{rsample} makes the path from \code{loss} back to the
encoder's parameters differentiable. The closed-form
Normal-Normal KL avoids MC noise.

\section{Constraints}
\label{sec:distributions:constraints}

Distribution parameters have constraints (variances are positive,
probabilities are in $[0, 1]$, simplex distributions sum to 1).
\Lucid\ exposes a \code{constraints} submodule that names these
constraints and a \code{transforms} submodule that bijects
between constrained and unconstrained spaces.

\subsection{Why constraints matter}
\label{ssec:distributions:why_constraints}

Optimising a positive parameter directly (e.g., a Gamma rate)
risks the optimiser producing a negative update that the
distribution then rejects with NaN. The standard remedy is to
parameterise via the unconstrained space and transform at the
distribution's input.

\begin{lstlisting}[style=lucid,language=LucidPython]
import lucid
from lucid.distributions import Gamma
from lucid.distributions.transforms import ExpTransform

# Parameterise log-rate as the unconstrained variable
log_rate = lucid.tensor(0.0, requires_grad=True)
rate = log_rate.exp()  # always positive

dist = Gamma(rate, 1.0)
# Optimise log_rate freely; rate stays in (0, infty).
\end{lstlisting}

\subsection{TransformedDistribution}
\label{ssec:distributions:transformed}

\code{TransformedDistribution(base, transform)} composes a base
distribution with a bijective transform. The
\code{log\_prob} correctly accounts for the Jacobian:
\[
\log p_Y(y) = \log p_X(\text{inv}(y)) - \log |J_{T}(\text{inv}(y))|.
\]

This is the mechanism for normalising flows: a chain of
transformed distributions, each transform learnable.

\section{Implementation Notes}
\label{sec:distributions:impl}

The distributions live in \code{lucid/distributions/}. Each
distribution is a Python class deriving from \code{Distribution}.

\subsection{Numpy-free}
\label{ssec:distributions:numpy_free}

Per Hard Rule H4, no distribution code imports numpy. The
implementations use only \Lucid\ primitives and special-function
calls. Erfinv-based Gaussian sampling on CPU,
\Lucid's Ziggurat on GPU.

\subsection{Composition with optim}
\label{ssec:distributions:compose_optim}

The distributions interact with the optimisers
(\chref{ch:optimizers}) through their parameters. A
\code{Distribution} created from \code{requires\_grad=True}
tensors backwards through both the sample and the log\_prob; the
optimiser updates the underlying tensors and the next sample
reflects the change.

\section{Univariate Continuous Distributions in Detail}
\label{sec:distributions:continuous_detail}

The most-used continuous distributions, with their densities,
log-likelihoods, and gradient identities.

\subsection{Normal}
\label{ssec:distributions:normal_detail}

Density:
\[
p(x; \mu, \sigma) \;=\; \frac{1}{\sigma\sqrt{2\pi}}
\exp\!\left(-\frac{(x - \mu)^2}{2\sigma^2}\right).
\]
Log-likelihood:
\[
\log p(x; \mu, \sigma) \;=\; -\tfrac{1}{2}\log(2\pi) - \log\sigma
- \frac{(x - \mu)^2}{2\sigma^2}.
\]
Gradients:
\[
\frac{\partial \log p}{\partial \mu} = \frac{x - \mu}{\sigma^2},
\qquad
\frac{\partial \log p}{\partial \sigma} = -\frac{1}{\sigma} +
\frac{(x - \mu)^2}{\sigma^3}.
\]
Entropy: $H = \tfrac{1}{2}\log(2\pi e \sigma^2)$.

\subsection{Gamma}
\label{ssec:distributions:gamma_detail}

Density (rate parameterisation, $\alpha$ shape, $\beta$ rate):
\[
p(x; \alpha, \beta) \;=\; \frac{\beta^\alpha}{\Gamma(\alpha)}
x^{\alpha - 1} e^{-\beta x},
\quad x > 0.
\]
Log-likelihood:
\[
\log p \;=\; \alpha \log \beta - \log\Gamma(\alpha) + (\alpha - 1)\log x - \beta x.
\]
Gradients (involving the digamma function $\psi$):
\[
\frac{\partial \log p}{\partial \alpha} = \log \beta - \psi(\alpha) + \log x,
\qquad
\frac{\partial \log p}{\partial \beta} = \frac{\alpha}{\beta} - x.
\]
Mean: $\alpha/\beta$. Variance: $\alpha/\beta^2$.

\subsection{Beta}
\label{ssec:distributions:beta_detail}

Density:
\[
p(x; \alpha, \beta) \;=\; \frac{1}{\mathrm{B}(\alpha, \beta)}
x^{\alpha - 1} (1 - x)^{\beta - 1},
\quad x \in (0, 1).
\]
Mean: $\alpha / (\alpha + \beta)$. The Beta is the conjugate
prior for the Bernoulli likelihood; this is the source of its
prominence in Bayesian inference.

\subsection{Dirichlet}
\label{ssec:distributions:dirichlet_detail}

The multivariate generalisation of Beta, defined on the
$(K-1)$-simplex:
\[
p(x; \alpha) \;=\; \frac{1}{\mathrm{B}(\alpha)}
\prod_{i=1}^K x_i^{\alpha_i - 1},
\quad x_i \geq 0,\ \sum_i x_i = 1,
\]
where $\mathrm{B}(\alpha) = \prod_i \Gamma(\alpha_i) /
\Gamma(\sum_i \alpha_i)$. Conjugate prior for categorical
likelihoods.

\subsection{Exponential and related}
\label{ssec:distributions:exp_related}

Exponential with rate $\lambda$:
\[
p(x; \lambda) \;=\; \lambda e^{-\lambda x},
\quad x \geq 0.
\]
Memoryless. Maximum-entropy distribution given a fixed mean.

Laplace (double exponential):
\[
p(x; \mu, b) \;=\; \frac{1}{2b} \exp\!\left(-\frac{|x - \mu|}{b}\right).
\]

Cauchy (heavy-tailed, infinite mean/variance):
\[
p(x; \mu, \gamma) \;=\; \frac{1}{\pi \gamma\left[1 + ((x - \mu)/\gamma)^2\right]}.
\]

\section{Discrete Distributions in Detail}
\label{sec:distributions:discrete_detail}

\subsection{Bernoulli}
\label{ssec:distributions:bernoulli}

Single trial with success probability $p$:
\[
P(X = 1) = p,\quad P(X = 0) = 1 - p.
\]
Log-likelihood:
\[
\log P(X = x) \;=\; x \log p + (1 - x) \log(1 - p).
\]
The sigmoid output of a logistic regression model is naturally
the parameter of a Bernoulli.

\subsection{Categorical}
\label{ssec:distributions:categorical}

The discrete distribution over $K$ classes with probability
vector $\pi$:
\[
P(X = k) = \pi_k,\quad k = 1, \ldots, K,\ \sum_k \pi_k = 1.
\]
The softmax output of a multi-class classifier is naturally
$\pi$. The cross-entropy loss is the negative log-likelihood:
\[
\mathcal{L} \;=\; -\log \pi_{y} \;=\; -y^\top \log \pi
\]
(where $y$ is the one-hot label).

\subsection{Binomial and Poisson}
\label{ssec:distributions:binomial_poisson}

Binomial: number of successes in $n$ Bernoulli trials:
\[
P(X = k) \;=\; \binom{n}{k} p^k (1 - p)^{n - k}.
\]
Poisson: count of events in fixed interval with mean rate
$\lambda$:
\[
P(X = k) \;=\; \frac{\lambda^k e^{-\lambda}}{k!}.
\]
Poisson is the limit of Binomial with $n \to \infty$, $p \to 0$,
$np = \lambda$.

\section{Multivariate Distributions}
\label{sec:distributions:multivariate}

\subsection{Multivariate Normal}
\label{ssec:distributions:mvn}

The density of $\mathcal{N}(\mu, \Sigma)$ on $\mathbb{R}^d$:
\[
p(x) \;=\; \frac{1}{(2\pi)^{d/2} |\Sigma|^{1/2}}
\exp\!\left(-\tfrac{1}{2}(x - \mu)^\top \Sigma^{-1} (x - \mu)\right).
\]
Log-likelihood (the form used computationally):
\[
\log p(x) \;=\; -\tfrac{d}{2}\log(2\pi) - \tfrac{1}{2}\log|\Sigma|
- \tfrac{1}{2}(x - \mu)^\top \Sigma^{-1} (x - \mu).
\]

\Lucid\ supports three parameterisations of $\Sigma$:
\begin{itemize}
  \item \code{covariance\_matrix}: full $\Sigma$.
  \item \code{precision\_matrix}: $\Sigma^{-1}$.
  \item \code{scale\_tril}: Cholesky factor $L$ with $\Sigma = L
    L^\top$.
\end{itemize}
The Cholesky parameterisation is the most numerically robust and
the default in most tutorials.

KL between two MVNs:
\[
D_{\mathrm{KL}}\!\bigl(\mathcal{N}(\mu_1, \Sigma_1) \,\|\, \mathcal{N}(\mu_2, \Sigma_2)\bigr)
= \tfrac{1}{2}\!\left[\operatorname{tr}(\Sigma_2^{-1} \Sigma_1)
+ (\mu_2 - \mu_1)^\top \Sigma_2^{-1} (\mu_2 - \mu_1)
- d + \log\frac{|\Sigma_2|}{|\Sigma_1|}\right].
\]

\subsection{Wishart and inverse Wishart}
\label{ssec:distributions:wishart}

The Wishart is the conjugate prior for the precision matrix of a
multivariate Normal in Bayesian inference. The inverse Wishart is
the prior for the covariance matrix.

\subsection{LKJ-Cholesky}
\label{ssec:distributions:lkj}

The LKJ distribution is a prior over correlation matrices. The
Cholesky parameterisation \code{LKJCholesky(d, concentration)}
returns the Cholesky factor of a correlation matrix.

\section{Variational Inference: A Worked Example}
\label{sec:distributions:vi_detail}

To illustrate the distribution machinery in context, we walk
through a variational autoencoder (VAE) loss derivation
\cite{kingma_vae_2014}.

\subsection{The setup}
\label{ssec:distributions:vi_setup}

A latent-variable model has the form $p(x, z) = p(x|z) p(z)$
with prior $p(z) = \mathcal{N}(0, I)$ and likelihood $p(x|z)$
parameterised by a neural network. The marginal likelihood
$p(x) = \int p(x|z) p(z)\, dz$ is intractable.

The variational approximation introduces $q(z|x)$, also
parameterised by a neural network (the encoder), and considers
the evidence lower bound (ELBO):
\[
\log p(x) \;\geq\; \mathbb{E}_{q(z|x)}[\log p(x|z)] - D_{\mathrm{KL}}(q(z|x) \,\|\, p(z)).
\]

\subsection{The ELBO terms}
\label{ssec:distributions:elbo_terms}

The reconstruction term $\mathbb{E}_{q(z|x)}[\log p(x|z)]$ is
estimated by sampling $z \sim q(z|x)$ (via reparameterisation
for differentiability) and computing $\log p(x | z)$.

The KL term has closed form for the case
$q(z|x) = \mathcal{N}(\mu(x), \mathrm{diag}\sigma^2(x))$ and
$p(z) = \mathcal{N}(0, I)$:
\[
D_{\mathrm{KL}}(q \,\|\, p) \;=\; \tfrac{1}{2} \sum_{i=1}^d
\bigl(\mu_i^2 + \sigma_i^2 - 1 - \log \sigma_i^2\bigr).
\]
This is the famous ``KL divergence to standard normal'' term that
appears in every VAE training loss.

\section{Score-Function Estimator}
\label{sec:distributions:score_estimator}

For distributions that do not admit reparameterisation (discrete
distributions, certain heavy-tailed continuous), the score-
function estimator (REINFORCE) is the alternative:
\[
\nabla_\theta \mathbb{E}_{x \sim p_\theta}[f(x)]
\;=\; \mathbb{E}_{x \sim p_\theta}[f(x) \nabla_\theta \log p_\theta(x)].
\]

\subsection{Variance}
\label{ssec:distributions:score_variance}

The score-function estimator is unbiased but has high variance.
The standard variance-reduction trick uses a baseline $b$ that
does not depend on $x$:
\[
\nabla_\theta \mathbb{E}[f(x)]
\;=\; \mathbb{E}[(f(x) - b) \nabla_\theta \log p_\theta(x)].
\]
The expectation is unchanged (since
$\mathbb{E}[\nabla \log p_\theta] = 0$), but the variance can be
reduced by choosing $b$ near $\mathbb{E}[f(x)]$.

This is the algorithmic basis of REINFORCE with a baseline, the
ancestor of modern policy-gradient algorithms in RL.

\subsection{Implementation pattern}
\label{ssec:distributions:score_pattern}

\begin{lstlisting}[style=lucid,language=LucidPython]
# Forward
dist = Categorical(logits=policy_logits)
action = dist.sample()
log_p = dist.log_prob(action)
reward = env.step(action)   # non-differentiable

# Score-function gradient
baseline = reward.mean().detach()
loss = -((reward - baseline) * log_p).mean()
loss.backward()
\end{lstlisting}

\section{Transformed Distributions and Normalising Flows}
\label{sec:distributions:flows}

\code{TransformedDistribution(base, transform)} composes a base
distribution with a bijective transform. By the change-of-variables
formula:
\[
\log p_Y(y) \;=\; \log p_X(T^{-1}(y)) - \log |J_T(T^{-1}(y))|,
\]
where $J_T$ is the Jacobian of $T$. \Lucid\ supports several
transforms:

\begin{itemize}
  \item \code{AffineTransform}: $y = a x + b$.
  \item \code{ExpTransform}: $y = e^x$.
  \item \code{SigmoidTransform}: $y = \sigma(x)$.
  \item \code{TanhTransform}: $y = \tanh(x)$.
  \item \code{StickBreakingTransform}: maps $\mathbb{R}^{K-1}$
    to the $K$-simplex.
\end{itemize}

\subsection{Normalising flows}
\label{ssec:distributions:flow_arch}

A \emph{normalising flow} chains several transforms:
\[
y = T_K \circ T_{K-1} \circ \cdots \circ T_1(x), \quad
x \sim p_X.
\]
The resulting density is
\[
\log p_Y(y) \;=\; \log p_X(x) - \sum_{k=1}^K \log |J_{T_k}|.
\]

The transforms are chosen so that both $T$ and $T^{-1}$ and
$\log |J|$ are tractable. Common choices: coupling layers (RealNVP),
autoregressive transforms (MAF, IAF), spline-based transforms
(neural spline flows).

\Lucid\ does not currently ship a flow library out of the box,
but the building blocks (\code{TransformedDistribution} plus the
transforms) make implementation straightforward.

\section{Constraint and Transform Machinery}
\label{sec:distributions:constraints_detail}

The \code{constraints} submodule names common parameter
constraints:
\begin{itemize}
  \item \code{positive}: $\theta > 0$.
  \item \code{interval(a, b)}: $a < \theta < b$.
  \item \code{simplex}: $\theta_i \geq 0$, $\sum_i \theta_i = 1$.
  \item \code{positive\_definite}: matrix is symmetric and PD.
  \item \code{lower\_cholesky}: lower-triangular with positive
    diagonal.
\end{itemize}

The \code{biject\_to(constraint)} function returns a transform
that maps unconstrained $\mathbb{R}^n$ to the constrained space:
\begin{itemize}
  \item For \code{positive}: \code{ExpTransform()}.
  \item For \code{interval(a, b)}: scaled sigmoid.
  \item For \code{simplex}: \code{StickBreakingTransform()}.
\end{itemize}

This pattern lets the user optimise an unconstrained variable
and have the constraint enforced by the transform.

\section{Summary and Forward References}
\label{sec:distributions:summary}

\code{lucid.distributions} provides 33 parameterised
probability-distribution classes with sample / log\_prob / entropy
/ KL methods. The reparameterisation trick supports
differentiable sampling for the 15 distributions that admit it;
the score-function estimator is the alternative for the rest.
The KL-divergence registry contains 20 closed-form pairs with a
Monte-Carlo fallback.

The next chapter, \chref{ch:einops}, closes
\partref{part:math_subsystems} with the tensor-rearrangement
algebra. \partref{part:neural_networks} then begins the
neural-network library proper.

\begin{lucidnote}
For the user: \code{lucid.distributions} is parity-compatible
with \refframework's \code{torch.distributions}. Most patterns
(rsample for differentiable, log\_prob for likelihood, kl\_divergence
for variational losses) work identically.
\end{lucidnote}
